% =============================================================================
% 2.5 Architecture du système
% =============================================================================
\section{Architecture du système}
\label{sec:ch2-architecture}

\subsection{Vue d'ensemble : deux horloges}

L'architecture s'organise autour de deux horloges indépendantes, représentées
à la figure~\ref{fig:deux-horloges}. Cette séparation n'est pas un détail
d'implémentation~: elle exprime que l'ingestion et l'exploitation n'ont ni
la même durée, ni le même contrat vis-à-vis de l'utilisateur.

\begin{figure}[htbp]
  \centering
  \resizebox{\linewidth}{!}{% Architecture générale : deux horloges indépendantes
\begin{tikzpicture}[node distance=4mm and 7mm]

  % --- Horloge A ---
  \node[btptitleband, text width=12.2cm, fill=btpblue] (ta)
    {Horloge A — Ingestion \textnormal{\small(asynchrone, non bloquante)}};
  \node[btpbox, text width=1.9cm, below=7mm of ta.south west, anchor=north west] (a1) {Dépôt};
  \node[btpbox, text width=1.9cm, right=of a1] (a2) {Stockage};
  \node[btpbox, text width=1.9cm, right=of a2] (a3) {Analyse\\structurelle};
  \node[btpbox, text width=1.9cm, right=of a3] (a4) {Découpage\\en extraits};
  \node[btpbox, text width=1.9cm, right=of a4] (a5) {Indexation};
  \foreach \i/\j in {1/2, 2/3, 3/4, 4/5} {\draw[btparrow] (a\i) -- (a\j);}
  \node[btplabel, below=2.5mm of a3, text width=8cm]
    {Traitement long, exécuté en arrière-plan par des travailleurs dédiés};

  % --- Horloge B ---
  \node[btptitleband, text width=12.2cm, fill=btpsky,
        below=16mm of ta.south west, anchor=north west] (tb)
    {Horloge B — Exploitation \textnormal{\small(synchrone, interactive)}};
  \node[btpboxalt, text width=1.9cm, below=7mm of tb.south west, anchor=north west] (b1) {Question};
  \node[btpboxalt, text width=1.9cm, right=of b1] (b2) {Recherche\\hybride};
  \node[btpboxalt, text width=1.9cm, right=of b2] (b3) {Contexte};
  \node[btpboxalt, text width=1.9cm, right=of b3] (b4) {Génération\\contrainte};
  \node[btpboxalt, text width=1.9cm, right=of b4] (b5) {Validation\\humaine};
  \foreach \i/\j in {1/2, 2/3, 3/4, 4/5} {\draw[btparrow] (b\i) -- (b\j);}
  \node[btplabel, below=2.5mm of b3, text width=8cm]
    {Traitement court, borné en temps, restituant un résultat vérifiable};

  % --- Liaison ---
  \draw[btparrowlight, dashed] (a5.south) -- ++(0,-0.55)
    node[btplabel, right, xshift=1mm, yshift=1mm] {base de connaissances partagée}
    -| (b2.north);

\end{tikzpicture}
}
  \caption{Architecture générale~: deux horloges indépendantes}
  \label{fig:deux-horloges}
\end{figure}

L'\btpterm{horloge A} capitalise. Elle est asynchrone, non bloquante, et
peut durer des dizaines de secondes. L'utilisateur dépose un fichier et
reçoit une confirmation immédiate~; le traitement se poursuit en
arrière-plan.

L'\btpterm{horloge B} exploite. Elle est synchrone, bornée en temps, et
doit restituer un résultat vérifiable en quelques secondes. Elle ne
s'exécute que sur des documents dont l'horloge A a confirmé l'indexation.

Les deux horloges partagent la base de connaissances. Elles ne partagent
pas le fil d'exécution~: une ingestion longue ne doit jamais geler une
recherche.

\subsection{Architecture en couches}

La figure~\ref{fig:architecture-couches} décline cette organisation en
cinq couches de responsabilité.

\begin{figure}[htbp]
  \centering
  \resizebox{0.92\linewidth}{!}{% Architecture en couches de la plateforme
\begin{tikzpicture}[node distance=4mm]

  \newcommand{\layerw}{12.4cm}

  \node[btptitleband, text width=\layerw, fill=btpdeep] (l1)
    {Couche présentation};
  \node[btpboxplain, text width=\layerw, below=1.5mm of l1, align=left, font=\scriptsize]
    (c1) {Interface web~: dépôt de documents, conversation, revue des sections
           générées, export des livrables};

  \node[btptitleband, text width=\layerw, fill=btpblue, below=of c1] (l2)
    {Couche service et orchestration};
  \node[btpboxplain, text width=\layerw, below=1.5mm of l2, align=left, font=\scriptsize]
    (c2) {Interface applicative HTTP, authentification et cloisonnement par
           organisation, routage des intentions, machine à états du workflow,
           points d'interruption pour la supervision humaine};

  \node[btptitleband, text width=\layerw, fill=btpsky, below=of c2] (l3)
    {Couche traitement};
  \node[btpboxplain, text width=\layerw, below=1.5mm of l3, align=left, font=\scriptsize]
    (c3) {Analyse structurelle des documents et reconnaissance optique,
           découpage hiérarchique, calcul des plongements, recherche hybride,
           ré-ordonnancement, extraction d'exigences, génération contrainte};

  \node[btptitleband, text width=\layerw, fill=btplight, text=btpdeep, below=of c3] (l4)
    {Couche persistance};
  \node[btpboxplain, text width=\layerw, below=1.5mm of l4, align=left, font=\scriptsize]
    (c4) {Base relationnelle (métadonnées, extraits, exigences, versions,
           décisions), index vectoriel (dense et creux), stockage objet
           (fichiers d'origine et exports), file de travaux};

  \node[btptitleband, text width=\layerw, fill=btpslate, below=of c4] (l5)
    {Couche infrastructure};
  \node[btpboxplain, text width=\layerw, below=1.5mm of l5, align=left, font=\scriptsize]
    (c5) {Conteneurisation des services, journalisation structurée,
           supervision des métriques, sauvegardes};

  \draw[btparrow, line width=1.4pt] ($(l1.west)+(-0.4,0)$) -- ($(c5.west)+(-0.4,0)$)
    node[midway, rotate=90, anchor=south, btplabel, yshift=1pt] {Descente des appels};

\end{tikzpicture}
}
  \caption{Architecture en couches de la plateforme}
  \label{fig:architecture-couches}
\end{figure}

Cette décomposition autorise le remplacement d'une brique --- le modèle de
plongement, le magasin vectoriel, le moteur d'orchestration --- sans
remettre en cause les contrats des couches adjacentes. Elle prépare aussi
la seconde phase~: un magasin graphe s'insère dans la couche de
persistance, un n\oe ud agentique dans la couche de traitement, sans
modifier l'interface ni le workflow de revue.

\subsection{Flux de données}

Trois flux principaux traversent l'architecture.

\paragraph{Flux d'ingestion.}
Fichier $\rightarrow$ stockage objet $\rightarrow$ analyse structurelle
$\rightarrow$ extraits $\rightarrow$ plongements $\rightarrow$ index
vectoriel et base relationnelle.

\paragraph{Flux de recherche.}
Question $\rightarrow$ plongement $\rightarrow$ recherche hybride filtrée
$\rightarrow$ ré-ordonnancement $\rightarrow$ contexte numéroté.

\paragraph{Flux de génération.}
Contexte et consigne $\rightarrow$ modèle de langage $\rightarrow$ section
structurée $\rightarrow$ validation déterministe $\rightarrow$ revue humaine
$\rightarrow$ export.

Le chapitre~\ref{chap:resultats} nomme les composants qui incarnent ces
flux. La méthodologie s'arrête aux rôles, afin que les choix de produits
restent substituables.
